\documentclass[10pt,a4paper]{article}
\usepackage[utf8]{inputenc}
\usepackage[T1]{fontenc}
\usepackage[spanish, es-tabla]{babel}
\usepackage{lmodern}
\usepackage{csquotes}
\usepackage{amsmath, amsfonts, amssymb, amsthm}
\usepackage{booktabs}
\usepackage{graphicx}
\usepackage{subcaption}
\usepackage{float}
\usepackage{ragged2e}
\usepackage{tabularx}
\usepackage{array}
\newcolumntype{Y}{>{\RaggedRight\arraybackslash}X}
\usepackage[table]{xcolor}
\usepackage[
    left=2.54cm,
    right=2.54cm,
    top=2.54cm,
    bottom=2.54cm
]{geometry}

\usepackage{xcolor}
\usepackage[most]{tcolorbox}
\usepackage{soul}

\usepackage{listings}
\usepackage{upquote}

\lstset{
    basicstyle=\footnotesize\ttfamily,
    frame=single,
    numbers=left,
    numberstyle=\tiny\color{gray},
    numbersep=8pt,
    inputencoding=utf8,
    extendedchars=true,
    breaklines=true,
    breakatwhitespace=false,
    showstringspaces=false,
    keepspaces=true,
    keywordstyle=\color{blue},
    commentstyle=\color{green!50!black},
    stringstyle=\color{red!70!black},
    tabsize=4
}
\usepackage{tikz}

\usetikzlibrary{
    arrows.meta,
    positioning,
    automata
}

\usepackage[
    colorlinks=true,
    linkcolor=blue!50!black,
    urlcolor=blue!60!black,
    citecolor=blue!50!black
]{hyperref}

%================================================
% CAJA DE EJERCICIOS
%================================================
\newtcbtheorem[number within=section]
{ejercicio}{Ejercicio}{
    enhanced,
    sharp corners,
    boxrule=0pt,
    leftrule=5pt,
    colframe=gray!30!black,
    colback=gray!5!white,
    coltitle=white,
    fonttitle=\bfseries,
    separator sign={.\ },
    before skip=10pt,
    after skip=10pt
}{ej}

%================================================
% CAJA DE TEOREMAS
%================================================
\newtcbtheorem[number within=section]
{mytheo}{Teorema}{
    enhanced,
    sharp corners,
    boxrule=0pt,
    leftrule=5pt,
    colframe=blue!40!black,
    colback=blue!5!white,
    coltitle=white,
    fonttitle=\bfseries,
    separator sign={.\ },
    before skip=10pt,
    after skip=10pt
}{th}

%================================================
% CAJA DE ALERTA
%================================================
\newtcolorbox{alerta}[1]{
    enhanced,
    sharp corners,
    boxrule=0pt,
    leftrule=5pt,
    colframe=red!40!black,
    colback=red!5!white,
    coltitle=white,
    fonttitle=\bfseries,
    title={#1},
    before skip=10pt,
    after skip=10pt
}

%================================================
% CAJA DE NOTAS
%================================================
\newtcolorbox{nota}[1]{
    enhanced,
    sharp corners,
    boxrule=0pt,
    leftrule=5pt,
    colframe=orange!40!black,
    colback=yellow!5!white,
    coltitle=white,
    fonttitle=\bfseries,
    title={#1},
    before skip=10pt,
    after skip=10pt
}


\begin{document}

\begin{titlepage}
    \centering

    %========================================
    % UNIVERSIDAD
    %========================================
    {\huge\textbf{Universidad Nacional Autónoma de México}}

    \vspace{0.3cm}
    \rule{\linewidth}{2mm}

    \vspace{0.3cm}
    {\large\textbf{Facultad de Estudios Superiores Acatlán}}

    \vspace{0.8cm}

    %========================================
    % LOGOTIPOS
    %========================================
    \noindent
    \begin{minipage}{0.45\textwidth}
        \centering
        \IfFileExists{../../../Archivos/Imagenes_Logos/acatlan.png}{
            \includegraphics[
                width=5cm,
                height=4cm,
                keepaspectratio
            ]{../../../Archivos/Imagenes_Logos/acatlan.png}
        }{
            \textit{Logotipo FES Acatlán}
        }
    \end{minipage}
    \hfill
    \begin{minipage}{0.45\textwidth}
        \centering
        \IfFileExists{../../../Archivos/Imagenes_Logos/SAT-BOT.jpg}{
            \includegraphics[
                width=5cm,
                height=4cm,
                keepaspectratio
            ]{../../../Archivos/Imagenes_Logos/SAT-BOT.jpg}
        }{
            \textit{Logotipo BOT}
        }
    \end{minipage}

    \vspace{0.8cm}

    %========================================
    % CARRERA Y TÍTULO
    %========================================
 

    \vspace{0.8cm}
    {\huge\textbf{Temas Selectos de Computación I}}

    \vspace{0.4cm}
    \rule{\linewidth}{1mm}

    \vspace{0.4cm}

    %========================================
    % PROYECTOS
    %========================================
    {\Large\textbf{Proyecto}}

    \vspace{0.3cm}
    {\large\textit{Análisis de Impacto y Planeación Arquitectónica: SAT-Bot Ciudadano 2026}}

    \vspace{6cm}

    %========================================
    % INTEGRANTES
    %========================================
    \rule{\linewidth}{0.5mm}

    \vspace{0.3cm}
    {\large\textbf{Integrantes}}

    \vspace{0.3cm}
    {\large Vazquez Perez Xochitl}

    \vspace{0.15cm}
    {\large Estrada Romero Meliza Edith}

    \vspace{0.15cm}
    {\large López Ramírez Ricardo}


    \vspace{0.5cm}
    \rule{\linewidth}{0.5mm}

    \vspace{0.5cm}

    %========================================
    % FECHA
    %========================================
    {\small Actualización: \today}

\end{titlepage}


%================================================
% ÍNDICE
%================================================
\pagenumbering{roman}

\tableofcontents

\clearpage

%================================================
% CONTENIDO
%================================================
\pagenumbering{arabic}
\setcounter{page}{1}

\section{Tipo de datos  a utilizar.}

El SAT-Bot procesará un flujo híbrido de datos para garantizar una respuesta
precisa y auditable:

\begin{itemize}
    \item Datos No Estructurados (Documentales): 
    Archivos PDF correspondientes a la Constancia de Situación Fiscal que 
    el ciudadano sube temporalmente a la plataforma.
    \item Datos Estructurados (Extraídos): 
    Entidades clave obtenidas mediante IA a partir de los PDFs, 
    tales como RFC, Razón Social, Régimen Fiscal, Código Postal y 
    Obligaciones Fiscales. Se manejarán en formato JSON.
    \item Datos Conversacionales (Procesamiento de Lenguaje Natural): 
    Historial de interacción entre el ciudadano y el modelo de lenguaje 
    (Vertex AI), incluyendo preguntas abiertas sobre dudas tributarias.
    \item Datos de Telemetría y Analítica: 
    Metadatos de uso del sistema 
    (latencia en milisegundos, códigos de estado HTTP, 200/500), 
    timestamps de las consultas y volumetría de uso para alimentar 
    los tableros de impacto.
\end{itemize}

\newpage    

\section{Elección de Base de Datos y Comparativa}

Se requiere una persistencia dual: 
una para el estado rápido de las conversaciones y otra para el análisis masivo de métricas.

\begin{table}[htbp]
    \centering
    \caption{Elección de Base de Datos y Comparativa}
    \label{tab:comparativa-bases-datos}

    \small
    \renewcommand{\arraystretch}{1.5}

    \begin{tabularx}{\textwidth}{
        p{2.5cm}
        p{3.0cm}
        Y
        Y
        Y
    }

        \toprule

        \rowcolor{blue!15}
        \textbf{Tipo de Base de Datos} &
        \textbf{Servicio GCP} &
        \textbf{Casos de Uso Ideales} &
        \textbf{Desventajas para este proyecto} &
        \textbf{Veredicto} \\

        \midrule

        \textbf{Relacional (SQL)} &
        Cloud SQL (PostgreSQL/MySQL) &
        Transacciones financieras, alta integridad relacional (ACID), cuentas de usuario complejas. &
        Requiere aprovisionar capacidad fija (instancias), lo que consume créditos constantemente aunque no haya tráfico. &
        \textbf{\textcolor{red}{Descartada}} temporalmente; sobrecualificada para guardar solo chats y métricas simples. \\

        \midrule

        \textbf{NoSQL (Documentos)} &
        Firestore &
        Guardar el estado de la sesión del chat, perfiles temporales, lecturas/escrituras de muy baja latencia y esquema flexible (JSON). &
        No es ideal para hacer consultas analíticas complejas o agregaciones masivas de datos históricos. &
        \textbf{\textcolor{green!50!black}{Seleccionada}} para la capa de la aplicación (Frontend/Backend), para guardar el historial del chat. \\

        \midrule

        \textbf{Data Warehouse} &
        BigQuery &
        Análisis masivo de datos, dashboards, consultas SQL sobre terabytes de información e integración nativa con Looker Studio. &
        No está diseñada para inserciones transaccionales de un solo registro en tiempo real para aplicaciones web. &
        \textbf{\textcolor{green!50!black}{Seleccionada}} para la capa analítica. Almacenará los logs de uso para el dashboard de impacto. \\

        \bottomrule

    \end{tabularx}
\end{table}

\vspace{2mm}

Utilizaremos Firestore como base de datos operativa para que Melissa conecte el contexto del chatbot, y 
BigQuery como motor analítico Serverless para que Xóchitl estructure los KPIs del proyecto sin necesidad de administrar servidores de bases de datos.

\newpage

\section{Análisis de Mercado y Oportunidad}

\vspace{2mm}

Soluciones actuales en el mercado


\vspace{2mm}

\begin{itemize}
    \item Orión (Chatbot Oficial del SAT): Es un bot basado en árboles de decisión rígidos (reglas).
    Carece de comprensión profunda de lenguaje natural, frecuentemente no resuelve 
    dudas específicas y obliga al usuario a navegar por menús predefinidos. 
    No permite subir documentos para su análisis.
    \item Heru / Taxfyle (Plataformas Contables Privadas): 
    Aplicaciones robustas (B2C) que conectan al usuario con contadores humanos o 
    automatizan declaraciones enteras. Son soluciones de pago por evento o suscripción.
    \item Bots de Despachos Contables (WhatsApp/Web): 
    Principalmente utilizados como embudos de ventas para captar clientes; 
    no ofrecen resolución técnica gratuita ni OCR automatizado de constancias.
\end{itemize}


\vspace{2mm}

Diferenciador y Oportunidad (Nuestra Propuesta de Valor):
SAT-Bot Ciudadano 2026 se posiciona como una herramienta de primer nivel de atención 
(Triage Fiscal) gratuita, de código abierto e impulsada por IA Generativa. 
A diferencia de Orión, nuestro bot permite la carga directa de documentos oficiales (PDF),
 utilizando Visión Artificial (Document AI) para "leer" la situación del ciudadano y 
 Vertex AI para traducir la jerga fiscal del SAT a un lenguaje empático, claro y 
 accesible para cualquier persona, resolviendo dudas inmediatas sin costo alguno.

\newpage

\section{Planteamiento de la Infraestructura en la Nube (Arquitectura GCP)}

\vspace{2mm}

La arquitectura ha sido diseñada bajo el paradigma Serverless (Sin Servidor),
lo cual garantiza que la aplicación pueda escalar a miles de usuarios si es necesario, pero escale a cero (costo cero) 
cuando no se utiliza, optimizando los créditos educativos.

\begin{itemize}
    \item Cómputo (Cloud Run):
    Elegido sobre Compute Engine (Máquinas Virtuales) porque nos permite desplegar 
    nuestro backend en Python/FastAPI mediante contenedores 
    Docker sin tener que parchar sistemas operativos ni gestionar redes. 
    Se cobra solo por los milisegundos de procesamiento.
    \item Almacenamiento (Cloud Storage): Actúa como una zona de aterrizaje (landing zone) temporal y segura.
    Elegido por su bajo costo y capacidad de disparar eventos. 
    Los PDFs de los ciudadanos se suben aquí e inmediatamente se borran tras su 
    análisis por seguridad de los datos personales.
    \item Inteligencia Artificial (Document AI, Vertex AI):
    Document AI (Form Parser): Elegido por ser un modelo preentrenado 
    para extraer llaves/valores de documentos formales, 
    superando ampliamente a un OCR tradicional (como Tesseract).

    \vspace{2mm}

    Vertex AI:Provee el cerebro conversacional LLM. 
    Elegido sobre la API pública de OpenAI por mantenerse dentro del 
    ecosistema GCP, garantizando que los datos fiscales de los ciudadanos
     no se utilicen para entrenar modelos públicos (privacidad empresarial).
    \item CI/CD (Cloud Build , Artifact Registry): 
    Automatizan el empaquetado del código en imágenes Docker seguras. 
    Elegidos por su integración nativa y fluida con Cloud Run.
\end{itemize}

\end{document}

